\chapterpage{Módulo 6 · Preparar la exposición}{Guion de demostración de 25 minutos}
Prepare una ruta principal con expedientes existentes y una alternativa de práctica en vivo. Así podrá explicar el sistema incluso si se agota la cuota del proveedor. No presente una consulta histórica como generación nueva.
\begin{tabularx}{\linewidth}{@{}p{18mm}YY@{}}
\toprule
\textbf{Minutos}&\textbf{Acción visible}&\textbf{Explicación sugerida}\\
\midrule
0--3&Presentar problema y conceptos.&«Convertimos una necesidad comercial en un diseño verificable. La IA asiste; la plataforma comprueba; una persona aprueba».\\
3--6&Explorador AW y WWI.&«Los metadatos describen fuentes diferentes. Una línea de pedido no es una línea de factura».\\
6--8&Catálogos por fuente.&«Las preguntas orientan, pero no inventan columnas. AW conserva sus personalizaciones; WWI tiene su propio catálogo».\\
8--12&Necesidad, consentimiento y viabilidad.&«Los derivables son candidatos. La estructura y los límites se revisan antes de generar». Usar capturas si no se hará una llamada.\\
12--16&Abrir \#98 y \#114.&«Estas propuestas conservan grano y documentos. El promedio por pedido/factura usa un conteo distinto».\\
16--19&Abrir expedientes \#14 y \#15.&«Aprobar no fue cargar. Aquí están la ejecución y la conciliación de origen y destino».\\
19--23&Panel, filtros y gráfico interactivo.&«El dato cambia con el filtro y conserva procedencia. El total de WWI no se supone USD».\\
23--25&Copiloto, exportación y límites.&«Contrasto indicador, evento, filtro y denominador. No afirmo causalidad ni compatibilidad no probada».\\
\bottomrule
\end{tabularx}
\sub{Lista de preparación del expositor}
\begin{itemize}
\item Probar acceso y permisos; confirmar ambas fuentes sin cambiar sus catálogos.
\item Localizar las propuestas y los expedientes antes de exponer. Si sus datos físicos ya no están, no prometer navegación en vivo del panel.
\item Preparar capturas y exportaciones con fecha, fuente y filtros. No mostrar claves.
\item Comprobar cuota sólo si se planean nuevas llamadas autorizadas. No prometer una latencia exacta.
\item Si se genera un caso nuevo, registrar su propio número. No reemplazar los identificadores de la evidencia histórica.
\end{itemize}
\textbf{Cierre sugerido:} «Este recorrido está acreditado en dos fuentes SQL Server. Rechazar una sugerencia sin respaldo y conservar evidencia es parte de la calidad del sistema».

\chapterpage{Módulo 6 · Aprendizaje activo}{Taller guiado para los tres perfiles}
Los tiempos son orientativos. La práctica de 80 minutos supone aplicación, fuentes, cuenta y proveedor preparados. El estudio de conceptos puede realizarse antes, sin consumir IA.
\begin{tabularx}{\linewidth}{@{}p{19mm}Yp{45mm}@{}}
\toprule
\textbf{Tiempo}&\textbf{Actividad}&\textbf{Producto del alumno}\\
\midrule
0--10&Leer conceptos, iniciar sesión y reconocer la fuente.&Explicar pedido, factura, fila y KPI con sus palabras.\\
10--20&Explorar ambas bases y una relación real.&Identificar PK, FK, fecha y grano candidato.\\
20--30&Comparar catálogos y escribir objetivo pequeño.&Texto, preguntas, periodicidad y límites.\\
30--45&Autorizar y validar; leer cada resultado.&Registro de directos, derivables y decisiones pendientes.\\
45--55&Generar si está autorizado y revisar contrato.&Grano, medida, fórmula, denominador y cobertura.\\
55--65&Aprobar y ejecutar un caso controlado, o consultar histórico si no alcanza el tiempo.&Identificar honestamente ejecución nueva o consulta.\\
65--75&Filtrar, abrir gráfico, preguntar y exportar.&Comparación de cifras y captura del alcance.\\
75--80&Redactar conclusión e incidente si corresponde.&Acta con esperado, obtenido, evidencia y límite.\\
\bottomrule
\end{tabularx}
\sub{Ruta por perfil}
\textbf{Usuario final / gerencia:} estudie conceptos y criterios de ventas; consulte panel, copiloto y exportaciones. Debe poder cuestionar un promedio sin administrar credenciales.

\textbf{Operador / analista autorizado:} añada catálogo, necesidad, diseño, personalización, aprobación, ETL y conciliación. Debe reconocer cuándo detener el proceso.

\textbf{Administrador / desarrollador:} añada entorno, seguridad, arquitectura, contratos, diagnóstico y pruebas. Debe resolver causas sin debilitar los controles ni confundir un problema de interfaz con uno de datos.
\begin{note}[Condición de aprobación del taller]
No se aprueba por llegar a un gráfico. El alumno debe identificar evento, grano, fuente, fórmula y límites, y saber dónde está la conciliación. Un caso negativo bien diagnosticado también demuestra aprendizaje; no debe disfrazarse de caso positivo.
\end{note}
\sub{Ejercicio de explicación entre compañeros}
Una persona hace de gerente y pregunta «¿puedo usar esto para decidir a qué cliente dar crédito?». La otra debe explicar que ventas o facturas por sí solas no prueban capacidad de pago: faltan datos y criterios adecuados. No invente una decisión crediticia desde un ranking comercial.

\chapterpage{Módulo 6 · Valores de referencia}{Controles para repetir sin confundir poblaciones}
Estos resultados se observaron en la auditoría del 3 de octubre de 2026 o en el expediente histórico señalado. \textbf{No son resultados nuevos producidos al crear este manual}. Identificadores y cifras dependen de la instalación y su corte.
\begin{tabularx}{\linewidth}{@{}Yrrr@{}}
\toprule
\textbf{Control}&\textbf{AW pedidos}&\textbf{WWI pedidos}&\textbf{WWI facturas}\\
\midrule
Propuesta / ejecución&98 / 14&96 / 13&114 / 15\\
Filas de detalle&121.317&231.412&228.265\\
Documentos distintos&31.465&73.595&70.510\\
Unidades&No se usa aquí&9.310.904&8.950.628\\
\bottomrule
\end{tabularx}
\sub{Importes y promedios que sí puede contrastar}
\begin{tabularx}{\linewidth}{@{}Yp{53mm}@{}}
\toprule
\textbf{Contexto}&\textbf{Valor observado}\\
\midrule
AW \#14: importe de todas las líneas&109.846.381,4250\\
AW \#14: promedio por pedido&3.491,07\\
AW \#14, 2014: pedidos / promedio&11.761 / 1.705,46\\
WWI \#13: importe según receta de pedidos&177.634.276,40\\
WWI \#15: importe facturado registrado&198.043.439,45\\
WWI \#15: promedio por factura&2.808,73\\
WWI \#15, 2015: total / facturas&62.090.220,81 / 22.250\\
WWI \#15, 2015: promedio por factura&2.790,57\\
\bottomrule
\end{tabularx}
\sub{Cuatro promedios de WWI que no deben confundirse}
En el expediente \#15: importe medio por línea aproximadamente \textbf{867,603178}; precio unitario medio simple \textbf{45,591688}; importe por unidad \textbf{22,126206}; importe por factura \textbf{2.808,728399}. La etiqueta, la composición del importe y el denominador hacen que respondan preguntas distintas.
\begin{note}[No es una comparación monetaria entre empresas]
AW mostró USD en el expediente observado; WWI conserva moneda de origen sin código único demostrado. Además, pedidos y facturas difieren. No calcule diferencias, participaciones o crecimiento entre las tres columnas como si fueran una sola serie de ventas comparable.
\end{note}
\textbf{Regla de contraste:} use el valor detallado y la precisión del expediente. Una tarjeta abreviada, como 109,8 M, no permite comprobar diferencias de centavos.

\chapterpage{Módulo 6 · QA para el alumno}{Pruebas positivas, negativas y registro de evidencia}
\begin{tabularx}{\linewidth}{@{}p{17mm}YY@{}}
\toprule
\textbf{Caso}&\textbf{Qué hacer}&\textbf{Resultado que debe comprobar}\\
\midrule
P-01&Consultar AW \#14 y filtrar un año presente.&Fuente/ejecución correctas; valores y gráficos responden al mismo filtro.\\
P-02&Consultar WWI \#15 y filtrar 2015.&Facturas, no pedidos; cifras coherentes con su expediente y denominador.\\
P-03&Volver de WWI a AW.&No quedan productos, territorios, revisión ni chat de la otra fuente como resultados propios.\\
N-01&Pedir clima externo con el texto siguiente.&Límite visible; no inventa clima ni sustituye meteorología por sensores internos.\\
N-02&Tras una revisión, editar el objetivo.&La revisión anterior no autoriza silenciosamente el alcance nuevo; revisar consentimiento.\\
N-03&Revisar una propuesta cuyo promedio sea de línea.&No aprobarlo como promedio documental. Corregir o etiquetar con honestidad.\\
N-04&Comparar factura con cobro en el copiloto.&No afirmar cobro sin evidencia; declarar limitación.\\
\bottomrule
\end{tabularx}
\begin{goal}
Analizar las ventas mensuales registradas y explicar cuánto influyó el clima de cada ciudad usando temperaturas externas que no están en esta base de datos.
\end{goal}
El negativo climático fue observado en NEED-20: se detectó el dato ausente y no se aceptó alcance parcial ni se generó ETL. El ensayo también encontró un problema temporal corregido antes del positivo final. \textbf{Un nuevo intento debe documentarse aparte}; no se ejecuta para demostrar que todo texto debe pasar.
\sub{Plantilla de ficha para la tesis o el taller}
\begin{tabularx}{\linewidth}{@{}p{42mm}Y@{}}
\toprule
Identificación&ID del caso, fecha, responsable y perfil utilizado.\\
Preparación&Fuente, instantánea, proveedor/modelo, permisos y versión del programa.\\
Entradas&Objetivo exacto, preguntas, periodicidad y límites aceptados.\\
Pasos&Acciones reproducibles; indicar consulta histórica o generación nueva.\\
Esperado / obtenido&Criterio antes de probar; observación real después, sin reescribir el esperado para hacerlo pasar.\\
Evidencia&Versión/ejecución si existen, capturas, filtros y conciliación.\\
Conclusión&Aprobado, fallido, parcial o no ejecutado; incidencia y riesgo pendiente.\\
\bottomrule
\end{tabularx}

\chapterpage{Módulo 6 · Qué hacer ante problemas}{Acceso, preparación e inteligencia artificial}
\sub{Primero deténgase, identifique y conserve}
No pulse varias veces una operación con efecto real. Registre pantalla, mensaje, fuente, versión, hora y última acción. No incluya credenciales. Distinga si falló el \textbf{acceso}, la \textbf{conexión}, el \textbf{proveedor}, el \textbf{contrato} o el \textbf{cálculo}.
{\small
\begin{tabularx}{\linewidth}{@{}p{29mm}YY@{}}
\toprule
\textbf{Síntoma}&\textbf{Acción segura del usuario}&\textbf{Cuándo escalar / qué evitar}\\
\midrule
No abre la página&Verificar dirección y si está en el equipo correcto; avisar al responsable técnico.&No borrar volúmenes ni reinstalar por intuición. Revisar servicios y puertos.\\
No puede entrar&Comprobar correo; usar el acceso asignado. Si venció sesión, volver a autenticar.&No compartir claves ni usar otra cuenta para saltarse permisos.\\
No aparece una opción&Confirmar perfil y permiso con administrador.&La vista ejecutiva no otorga permisos. No modificar roles para una demostración sin autorización.\\
Fuente no lista&Revisar conexión probada, habilitación e instantánea.&El administrador revisa acceso de sólo lectura, servidor y base. No duplicar conexión por cada fallo.\\
No hay metadatos&Crear primera instantánea con permiso. Si hay una, revisar contexto y versión.&Actualizar estructura no carga ventas. No cambiar metadatos para que acepte una referencia falsa.\\
Validar deshabilitado&Objetivo entre 20 y 2.000 caracteres, pregunta, período, destino y consentimiento.&Leer mensaje; adoptar una sugerencia obliga a revisar otra vez el contexto.\\
Sugerencia bloqueada&Mantener texto propio; revisar qué referencia o capacidad falta.&No asumir que la necesidad comercial es inválida sólo porque una reformulación falló.\\
Cuota / límite temporal&Conservar objetivo; esperar el intervalo indicado o usar históricos.&No reintentar sin límite ni cambiar proveedor sin coordinar la configuración activa.\\
JSON inválido o truncado&No usar la respuesta descartada. Reintentar de forma acotada cuando la aplicación lo permita.&Desarrollador revisa categoría del error y contrato. No copiar JSON parcial para aprobarlo manualmente.\\
Necesidad no disponible&Reducir alcance con decisión explícita o aportar una fuente autorizada.&No inventar datos. Si lo ausente es central para la decisión, no continúe como si estuviera resuelto.\\
\bottomrule
\end{tabularx}}
\begin{note}[Durante una exposición]
Diga qué falló, muestre que se conservó el trabajo y pase al expediente de referencia identificado como histórico. Esa recuperación es demostrable; ocultar el fallo y afirmar que la generación fue exitosa no lo es.
\end{note}

\chapterpage{Módulo 6 · Qué hacer ante problemas}{Diseño, ETL, panel y coherencia comercial}
{\small
\begin{tabularx}{\linewidth}{@{}p{29mm}YY@{}}
\toprule
\textbf{Síntoma}&\textbf{Qué comprobar y hacer}&\textbf{Qué no debe hacer}\\
\midrule
Error bloqueante en propuesta&Leer referencias, grano, tipos y cobertura. Personalizar una opción válida o generar una nueva versión corregida.&No aprobar ni eliminar la regla para desbloquear.\\
Histórico con advertencias&Distinguir aprobación histórica, compatibilidad actual y datos disponibles. Verificar evidencia y registrar diferencias.&No retirar aprobación ni descartar versiones sólo por consultarlas.\\
No puede crear otra necesidad&Usar Crear nueva propuesta desde el histórico.&No cerrar sesión como método normal ni borrar el resultado anterior.\\
KPI Sin valor&Revisar estado, fórmula, medida base, denominador y filtro.&No sustituir nulo por cero ni afirmar «no hubo ventas» sin respaldo.\\
Promedio incoherente&Identificar si es por línea, documento o unidad. Comparar totales del mismo alcance.&No cambiar sólo la etiqueta para justificar otra fórmula.\\
ETL fallido&Abrir expediente, leer etapa y error; técnico revisa registros y conectividad.&No ejecutar repetidamente ni eliminar datamarts anteriores.\\
Diferencia en conciliación&Conservar resultados, revisar población, uniones, precisión y fórmula antes de publicar.&No normalizar la diferencia como aceptable sin criterio documentado.\\
Etiqueta pendiente&Revisar mapeos y usar publicación o reintento de interpretación sin repetir ETL.&No traducir identificadores o nombres personales ni cargar de nuevo por un texto.\\
Moneda desconocida&Usar Comprobar divisa sin repetir ETL si aparece; mantener moneda de origen si no se prueba.&No asumir USD ni convertir importes sin evidencia.\\
Panel ausente&Comprobar fuente, permiso, ejecución conciliada/publicada y tablas físicas.&Un historial guardado no garantiza datos físicos consultables.\\
Gráfico y chat distintos&Comparar evento, indicador, filtros, Top N, denominador y redondeo.&No asumir cruce de fuentes antes de revisar el alcance interpretado; si hay cruce real, detener análisis.\\
Fuente cruzada visible&Registrar pasos/captura y detener conclusiones; escalar a QA/desarrollo.&No ocultar el problema cambiando etiquetas ni reconstruir a ciegas ambas fuentes.\\
\bottomrule
\end{tabularx}}
\sub{Paquete mínimo para solicitar ayuda}
Pantalla y acción; fuente/instantánea; propuesta/ejecución; objetivo sin secretos; indicador y filtros; esperado y obtenido; fecha; captura legible. El técnico puede añadir categoría de error y registros saneados. Un caso reproducible vale más que «no funciona».

\chapterpage{Módulo 6 · Evaluación de aprendizaje}{Preguntas para comprobar que comprendió BI + IA}
Intente responder antes de leer la siguiente página. No hace falta memorizar código.
\begin{enumerate}
\item ¿En qué se diferencia una fuente OLTP de un datamart?
\item ¿Qué es el grano? Si una factura tiene tres líneas, ¿cuántas facturas cuenta?
\item ¿Qué diferencia existe entre dimensión Producto y atributo Categoría?
\item ¿Qué demuestra una FK y qué no demuestra sobre los importes?
\item ¿Por qué Mensual no significa perder el detalle ni programar ETL mensual?
\item ¿Pedido, factura y cobro pueden usarse como sinónimos? Dé un contraejemplo.
\item ¿Cuándo no debe llamar venta neta al importe de una columna?
\item ¿Cuál es la diferencia entre promedio de línea, por documento y por unidad?
\item ¿Por qué el producto con la línea más cara puede no ser el producto líder?
\item ¿Qué diferencia existe entre costo estándar, último costo y costo histórico?
\item ¿Qué hace la IA al sugerir una necesidad y qué sigue decidiendo la persona?
\item ¿Qué significa derivable? ¿Ya está conciliado ese cálculo?
\item ¿Aprobar una propuesta ejecuta ETL? ¿Verificar evidencia lo ejecuta?
\item ¿Los Hallazgos explicables del panel necesitan una llamada al LLM?
\item Si una respuesta del copiloto usa otro año, ¿debe coincidir con el panel actual?
\item ¿Qué significa que las filas concilien? ¿Certifica toda la interpretación comercial?
\item ¿Qué tres datos de procedencia debe mirar antes de leer el tablero?
\item ¿Por qué «moneda de origen» puede ser más correcto que mostrar USD?
\item ¿Cambiar Vista ejecutiva a Vista analítica demuestra que otro rol tiene acceso?
\item ¿Qué debe registrar cuando una prueba no puede completarse por cuota?
\end{enumerate}
\begin{good}[Criterio de dominio]
Puede usar la herramienta con criterio cuando responde estas preguntas con ejemplos, identifica la evidencia de una cifra y sabe detener una aprobación. Saber navegar sin comprender el evento y el denominador todavía no es dominar BI.
\end{good}

\chapterpage{Módulo 6 · Respuestas de estudio}{Respuestas comentadas y errores que evitan}
{\small
\begin{enumerate}
\item La fuente registra operaciones; el datamart organiza un alcance analítico. Aquí SQL Server es fuente y PostgreSQL almacena la materialización.
\item El grano define una fila. Tres líneas de una factura son tres hechos de detalle, pero un documento distinto. Evita inflar conteos.
\item Producto es el conjunto descriptivo; Categoría es uno de sus atributos. No es necesario crear una dimensión independiente por cada campo.
\item Una FK acredita una relación declarada entre claves. No acredita pago, composición fiscal ni calidad económica de los valores.
\item Mensual agrupa por año y mes. La fila puede seguir siendo una línea; la agenda de cargas es otro asunto.
\item No. Un pedido puede no haberse facturado; una factura puede estar pendiente de cobro. Cada evento requiere fuente y fecha propias.
\item Cuando no se conoce si incluye descuentos, impuestos, ajustes o qué población representa. Use importe registrado y documente límites.
\item Línea: importe / filas; documento: importe / documentos distintos; unidad: importe / unidades. El filtro y la base del importe deben coincidir.
\item Una línea máxima es una observación individual; el liderazgo por producto depende de sumar sus líneas. Confundirlos cambia la decisión.
\item Estándar es una referencia; último costo describe el valor disponible más reciente; histórico corresponde al momento/evento si existe evidencia. No son intercambiables.
\item La IA propone una redacción según contexto. La persona elige objetivo, acepta límites y aprueba después; la plataforma comprueba la salida.
\item Existe una ruta o receta candidata. Debe implementarse en el contrato y luego ejecutarse y conciliarse. No equivale a dato ya calculado.
\item No en ambos casos. Aprobar habilita una propuesta elegible; verificar comprueba estructura. La carga requiere la confirmación específica de materialización.
\item No: esos resúmenes se calculan por reglas. La conversación del copiloto sí puede interpretar y redactar mediante IA.
\item Sólo si usa el mismo alcance. Lea consulta interpretada y filtros; una consulta permitida no necesariamente cambia el tablero.
\item En estos ensayos acredita igualdad del detalle cargado; se revisan también medidas/KPI. No certifica por sí sola historia de costos, cobro o interpretación fiscal.
\item Fuente, propuesta y ejecución. Luego revise indicador, unidad, filtros y fecha de la carga.
\item Porque puede no existir evidencia de una moneda ISO única. Inventar un símbolo añade una certeza falsa.
\item No. Son modos de presentación. Debe probarse una cuenta autenticada con sus permisos reales.
\item Entradas, proveedor/modelo, hora, error, estado alcanzado y alternativa utilizada. Marque parcial/no ejecutado; no lo cuente como éxito de generación.
\end{enumerate}}
\sub{Ejercicio final de síntesis}
Explique en dos minutos una cifra del panel con esta estructura: \emph{«Proviene de esta fuente y ejecución; mide este evento con este grano; usa esta fórmula y estos filtros; se concilió así; y no permite concluir esto otro»}. Si puede hacerlo, está conectando operación, BI y juicio comercial.

\chapterpage{Referencia y trazabilidad documental}{Fuentes, vigencia y uso académico}
\sub{Qué se utilizó para construir esta guía}
\begin{enumerate}[label=R\arabic*.]
\item \textbf{Auditoría operativa BI + IA, 3 de octubre de 2026.} Archivo del proyecto \path{docs/testing/auditoria-operativa-2026-10-03.md}. Casos QA-01 a QA-16 y ampliación NEED-01 a NEED-21. Sustenta resultados, incidentes, alcance y limitaciones.
\item \textbf{Manual técnico del prototipo.} \path{docs/manual-tecnico/manual-tecnico.tex}. Secciones de operación, asistencia, materialización, analítica, configuración y recuperación. Se contrastaron sus instrucciones con el código actual.
\item \textbf{Arquitectura, alcance y réplica.} \path{docs/architecture.md}, \path{docs/scope.md}, \path{docs/replication-guide.md} y \path{AGENTS.md}. Sustentan la separación de responsabilidades y el trabajo de desarrollo.
\item \textbf{Interfaz vigente.} \path{frontend/src/App.tsx} y \path{frontend/src/AnalyticsPage.tsx}. Sustentan nombres de botones, estados, rutas, consentimiento y filtros; algunas capturas históricas muestran una disposición anterior.
\item \textbf{Contratos y servicios.} Módulos backend de metadatos, parámetros/proveedores, copiloto, ETL, analítica y reportes; Makefile y flujos CI/CD. Se revisaron para no atribuir al LLM tareas determinísticas ni prometer motores no implementados.
\item \textbf{Registros persistidos consultados en sólo lectura.} Catálogos, propuestas \#96/\#98/\#114 y expedientes \#13/\#14/\#15 del entorno de demostración. Sustentan la diferenciación de eventos y los controles de agregados.
\end{enumerate}
\sub{Procedencia de las capturas}
Las doce imágenes incluidas provienen de \path{docs/testing/evidencias/qa-20261003/}. Sus nombres originales se conservan en la carpeta \path{figuras/} del paquete LaTeX. Se reutilizan como evidencia real de aquella auditoría; no son interfaces recreadas ni pruebas nuevas realizadas durante la redacción.

Las capturas cubren explorador, catálogo, consentimiento, viabilidad, interpretación IA, aprobación, etapas de ETL, conciliación, panel, conversación y distinción de pedidos. La guía explica qué demuestra cada imagen y qué no.
\sub{Qué sigue pendiente fuera de este documento}
Una ejecución nueva de cada texto propuesto; pruebas de cuentas autenticadas de los roles nuevos; validación integral de otros motores/modelos/esquemas; y cualquier afirmación histórica o financiera que requiera datos adicionales. Esta guía no transforma esos pendientes en resultados aprobados.
\begin{good}[Cómo citarla en su trabajo]
Puede identificarla como \emph{Guía de formación desde cero: BI + IA aplicada a ventas}, versión 1.0, 3 de octubre de 2026, material de apoyo del prototipo. Para afirmar un resultado experimental, cite además el caso de auditoría y su evidencia; no use una ficha propuesta como si fuera una observación obtenida.
\end{good}
